% Standalone problem document, generated from the adjacent README.md.
% Compile with XeLaTeX. Mathematical content is maintained in Markdown.
\documentclass[11pt,a4paper]{article}
\usepackage[fontsize=10.5pt]{scrextend}
\usepackage[margin=22mm,headheight=15pt,headsep=9mm,footskip=11mm]{geometry}
\usepackage{fontspec}
\setmainfont{texgyrepagella}[Extension=.otf,UprightFont=*-regular,BoldFont=*-bold,ItalicFont=*-italic,BoldItalicFont=*-bolditalic]
\setsansfont{texgyreheros}[Extension=.otf,UprightFont=*-regular,BoldFont=*-bold,ItalicFont=*-italic,BoldItalicFont=*-bolditalic]
\setmonofont{texgyrecursor}[Extension=.otf,UprightFont=*-regular,BoldFont=*-bold,ItalicFont=*-italic,BoldItalicFont=*-bolditalic,Scale=MatchLowercase]
\usepackage{amsmath,amssymb}
\usepackage{unicode-math}
\setmathfont{texgyrepagella-math.otf}
\usepackage{xcolor,microtype,enumitem,needspace,fancyhdr}
\usepackage{bookmark}
\definecolor{ink}{HTML}{17344B}
\definecolor{accent}{HTML}{197D80}
\definecolor{muted}{HTML}{596773}
\hypersetup{colorlinks=true,linkcolor=accent,urlcolor=accent,pdftitle={MI-09 - Sharp
Schatten triangle constants for the arithmetic symmetric
modulus},pdfauthor={Open Problems in Numerical Linear Algebra}}
\urlstyle{same}
\setlength{\parindent}{0pt}
\setlength{\parskip}{5pt plus 1pt minus 1pt}
\linespread{1.04}
\setlength{\emergencystretch}{3em}
\widowpenalty=10000
\clubpenalty=10000
\displaywidowpenalty=10000
\allowdisplaybreaks[1]
\setlist{nosep,leftmargin=1.5em}
\providecommand{\tightlist}{\setlength{\itemsep}{0pt}\setlength{\parskip}{0pt}}
\providecommand{\pandocbounded}[1]{#1}
\pagestyle{fancy}
\fancyhf{}
\fancyhead[L]{\sffamily\footnotesize\color{muted}OPEN PROBLEMS IN NUMERICAL LINEAR ALGEBRA}
\fancyhead[R]{\sffamily\bfseries\footnotesize\color{accent}MI-09}
\fancyfoot[L]{\parbox[b]{0.9\textwidth}{\sffamily\fontsize{8}{10}\selectfont\color{muted}Editorial ratings; a bounded search cannot certify that a problem remains unsolved.\\Literature check: 2026-09-08}}
\fancyfoot[R]{\sffamily\footnotesize\color{muted}\thepage}
\renewcommand{\headrulewidth}{0.3pt}
\renewcommand{\headrule}{\hbox to\headwidth{\color{accent}\leaders\hrule height \headrulewidth\hfill}}
\makeatletter
\renewcommand{\section}{\Needspace{5\baselineskip}\@startsection{section}{1}{0pt}{12pt plus 2pt minus 2pt}{4pt}{\sffamily\bfseries\large\color{ink}}}
\renewcommand{\subsection}{\Needspace{4\baselineskip}\@startsection{subsection}{2}{0pt}{9pt plus 1pt minus 1pt}{3pt}{\sffamily\bfseries\normalsize\color{ink}}}
\makeatother
\setcounter{secnumdepth}{0}
\begin{document}
{\sffamily\small\color{muted}Matrix inequalities and norms\par}
\vspace{3pt}
{\raggedright\hyphenpenalty=10000\exhyphenpenalty=10000\sffamily\bfseries\fontsize{21}{24}\selectfont\color{ink}Sharp
Schatten triangle constants for the arithmetic symmetric modulus\par}
\vspace{7pt}
{\sffamily\small\textbf{Difficulty:} challenging\quad\textbf{Importance:} interesting
to specialist\par}
\vspace{4pt}
{\color{accent}\hrule height 0.8pt}
\vspace{6pt}
\subsection{Problem statement}\label{problem-statement}

For \(X\in\mathbb C^{n\times n}\) set
\(S(X)=((X^*X)^{1/2}+(XX^*)^{1/2})/2\). Let
\(\|X\|_p=(\sum_i\sigma_i(X)^p)^{1/p}\) for \(1\le p<\infty\), with
\(\|X\|_\infty=\sigma_1(X)\). Determine, for all integers \(m,n\ge2\)
and all \(p\in[1,\infty]\), the exact value

\[c_p^{\rm sym}(m,n)=\sup_{(A_1,\ldots,A_m)\ne(0,\ldots,0)}
\frac{\|S(A_1+\cdots+A_m)\|_p}{\|S(A_1)+\cdots+S(A_m)\|_p},
\qquad A_j\in\mathbb C^{n\times n}.\]

The denominator is positive on the stated domain. A resolution should
identify the value for the unsolved parameters while retaining known
endpoint values.

\subsection{Why it matters}\label{why-it-matters}

This asks how much the symmetric positive representative of a matrix sum
can exceed the sum of the individual representatives in Schatten norm.
Sharp constants quantify the loss in this way of bounding nonnormal
matrix sums.

\subsection{References}\label{references}

\begin{enumerate}
\def\labelenumi{\arabic{enumi}.}
\tightlist
\item
  T. Zhang, \emph{Operator symmetric moduli and sharp triangle
  inequalities}, arXiv:2603.01046v1 (1 March 2026), Problem 1.17 and §8.
  \href{https://arxiv.org/html/2603.01046}{Primary text}.
\item
  J.-C. Bourin and E.-Y. Lee, \emph{Some hybrid matrix triangle
  inequalities}, arXiv:2606.29188v1 (28 June 2026), introduction and
  Theorem 1.3, for the dimension-independent majorization bound.
  \href{https://arxiv.org/html/2606.29188}{Primary text}.
\end{enumerate}

\subsection{Status check --- 2026-09-08}\label{status-check-2026-09-08}

The source proves \(c_1^{\rm sym}=1\) and \(c_\infty^{\rm sym}=\sqrt2\)
when \(n\ge3\), and gives bounds for the other cases. Its latest version
remains v1. Searches included
\textit{arithmetic symmetric modulus optimal Schatten constants},
\textit{Zhang Problem 1.17 symmetric modulus}, and both paper titles
with \textit{2026 sharp}. No general determination was found. Problem
1.11's two-dimensional restriction is subsumed here and is not counted
as an additional entry. Unlike the unitary-orbit conjecture, this
problem asks for exact norm extrema, not a positive-semidefinite
domination.
\end{document}
